%!TEX root = ../main.tex
\chapter{Metrics}
\label{ch:metrics}

Both training regimes used in this work---imitation of recorded flights
(Section~\ref{sec:training-supervised}) and reinforcement learning in
simulation (Section~\ref{sec:training-rl})---produce a single headline number
that the optimiser minimises: a loss in the first case, a negative expected
return in the second. Neither number answers the question that actually
matters, which is whether the resulting network flies the robot well. A
supervised loss can fall while the network learns nothing but ``predict zero'';
an episodic return can rise while the vehicle learns to crash less rather than
to reach its target. This chapter defines every quantity that is logged
alongside the headline number, states what each one measures, and explains why
it was added. The metrics are grouped by training regime, since the two ask
different questions of the model, and each group ends with a summary table.

Every metric named here is a column in a CSV file written during training: one
row per epoch for the supervised trainer, one row per iteration for the
reinforcement trainer. The column names are quoted verbatim in monospace, so
that a figure in Chapter~\ref{ch:training} can always be traced back to its
definition here.

\section{Conventions and Notation}
\label{sec:metrics-notation}

The network's output is a velocity command
$\mathbf{y} = (v_x, v_y, v_z, \omega_x, \omega_y, \omega_z) \in \mathbb{R}^{6}$:
three linear velocities in \si{\metre\per\second} and three angular velocities
in \si{\radian\per\second}. Components are indexed $i = 0,\dots,5$ in the order
above, matching the \met{_dim_<i>} suffix of the log columns. The demonstrated
(target) command is $\mathbf{y}$ and the network's prediction is
$\hat{\mathbf{y}}$; $N$ denotes the number of samples over which an average is
taken, which in sequence mode counts every time step of every window.

Each robot in the fleet declares its own kinematic limits inside the
\texttt{state} block of the observation. The per-component \emph{velocity cap}
$c_i > 0$ is the largest magnitude the platform accepts on component $i$;
a cap of zero declares that the axis does not exist on that robot (for example
vertical velocity on a ground vehicle).

For the reinforcement learning metrics the following terms are used with their
usual meaning, defined formally in Section~\ref{sec:rppo-gae}. An
\emph{episode} is one attempt at the navigation task, from spawn to a terminal
condition or to the time limit. A \emph{rollout} is a fixed number of
environment steps collected from every parallel environment with the current
policy; an \emph{iteration} is one rollout followed by one policy update, and
a \emph{mini-batch} is a subset of the rollout the update is computed on.
$\pi_\theta(a \mid s)$ is the stochastic policy, $V(s)$ the critic's estimate of
the state value, $\gamma$ the discount factor, $A_t$ the advantage of the action
taken at step $t$ and $R_t$ the corresponding bootstrapped return.

\section{Supervised Training Metrics}
\label{sec:metrics-supervised}

The supervised trainer learns to reproduce, at each instant of a recorded
flight, the command the human pilot or the autonomous pilot gave from the observation the robot had.
The questions its metrics answer are, in order: is the model learning anything
beyond a trivial rule; \emph{what} kind of behaviour is it reproducing; and is
the machinery (encoder, recurrent cells, gradients) healthy while it does so.
Every metric below is computed per batch and averaged over the epoch, and is
reported for the training split (prefix \met{train_}) and for the validation
split (prefix \met{val_}). A third family, \met{val_cal_}, is described in
Section~\ref{sec:metrics-calibrated}.

\subsection{The Training Loss: Cap-Scaled Mean Squared Error}
\label{sec:metrics-loss}

The loss that is optimised, and the quantity reported as \met{train_loss} and
\met{val_loss}, is a mean squared error in which every component of the command
is first divided by the robot's declared velocity cap:
%
\begin{equation}
  \mathcal{L}
  = \frac{1}{N}\sum_{n=1}^{N}\;\frac{1}{6}\sum_{i=0}^{5}
    \left(\frac{\hat{y}_{n,i} - y_{n,i}}{s_{n,i}}\right)^{2},
  \qquad
  s_{n,i} =
  \begin{cases}
    c_{n,i} & \text{if } c_{n,i} > \varepsilon,\\
    1       & \text{otherwise},
  \end{cases}
  \label{eq:capscaled}
\end{equation}
%
with $\varepsilon = 10^{-3}$ and $c_{n,i}$ the cap declared by the robot that
recorded sample $n$.

A plain MSE on raw commands is dominated by whichever component carries the
most energy. On the corpus used here the six components differ by orders of
magnitude in variance, so the gradient reaching the small ones is negligible
and they are effectively never learnt. The textbook remedy---dividing by the
per-component standard deviation of the dataset---is inappropriate for two
reasons specific to this problem. First, components can be identically zero
across the whole corpus, so their scale would be an arbitrary floor. Second, the
scale would shift every time new recordings were added, silently re-weighting
the loss between runs that are meant to be comparable.

Dividing by the robot's cap instead makes the error a \emph{fraction of what
the robot can do}. A ground robot capped at \SI{0.3}{\radian\per\second} and a
drone capped at \SI{1.0}{\radian\per\second} are then directly comparable, the
scale is a property of the platform rather than of the corpus, and a
heterogeneous fleet needs no special handling because every sample carries its
own caps. A cap below $\varepsilon$ is treated as ``this axis does not exist'':
the component is left in raw units, its target is necessarily zero, and the
loss simply drives the prediction to zero there without inflating its weight.

Two auxiliary values accompany the loss. The learning rate in effect at the
end of the epoch is logged as \met{lr}, so that a change of slope in the loss
curve can be matched to the schedule. When step-level logging is enabled,
\met{quick_val_loss} is the same loss evaluated on a handful of validation
batches part-way through an epoch, a cheap generalisation check for epochs that
last hours.

\subsection{Reference Predictors and Skill Scores}
\label{sec:metrics-r2}

A loss value on its own has no scale: \num{0.01} may be excellent or useless
depending on how large the commands are. The trainer therefore evaluates two
zero-parameter reference predictors on every batch and expresses the model's
error relative to them.

\paragraph{The zero predictor.}
\met{baseline_zero_loss} is Equation~\eqref{eq:capscaled} evaluated with
$\hat{\mathbf{y}} = \mathbf{0}$: the loss of a model that always commands
``hold still''. A run whose loss sits at or below this line has learnt nothing
beyond predicting nothing, which is the signature of the most common failure
mode observed---a loss that drops once in the first epoch and never moves
again.

The same comparison is made per component, with a plain (unscaled) squared
error so that the yardstick does not depend on the choice of loss:
%
\begin{equation}
  \text{\met{mse_dim_i}} = \frac{1}{N}\sum_{n}(\hat{y}_{n,i} - y_{n,i})^{2},
  \qquad
  \text{\met{baseline_dim_i}} = \frac{1}{N}\sum_{n} y_{n,i}^{2},
\end{equation}
%
and their ratio gives a coefficient of determination against the zero
predictor,
%
\begin{equation}
  \text{\met{r2_dim_i}} = 1 - \frac{\text{\met{mse_dim_i}}}{\text{\met{baseline_dim_i}}}.
  \label{eq:r2}
\end{equation}
%
A value of $1$ is a perfect fit, $0$ means the model is no better than
predicting zero, and a negative value means it is worse. The reason for
reporting this \emph{per dimension} is that on this corpus roughly \SI{78}{\percent}
of the target energy is concentrated in a single component of the six, so an
aggregate error that improves says almost nothing about the other five; a
model can score well overall while having learnt only the dominant one.

Two details of the computation matter. The ratio in
Equation~\eqref{eq:r2} is formed \emph{once per epoch} from the two epoch
means, not per batch and then averaged, because the pooled $R^2$ is a ratio of
means and not a mean of ratios. And when \met{baseline_dim_i} is exactly zero
the split never commands that component; the ratio is undefined, the column is
omitted rather than filled with an infinity, and the raw \met{mse_dim_i} is
kept, where it reads as pure error.

\paragraph{The persistence predictor.}
The observation already contains the robot's current linear and angular
velocity, which is the same physical quantity as the command one step later
and strongly correlated with it. Echoing that velocity is therefore a
zero-parameter rule that is far harder to beat than silence, and a model that
does not beat it is not worth its parameters. The loader exposes this echo as
$\mathbf{p}_n$ (masked to zero on the components the robot cannot reach or the
corpus never commands), and the trainer logs
%
\begin{equation}
  \text{\met{persist_dim_i}} = \frac{1}{N}\sum_{n}(p_{n,i} - y_{n,i})^{2},
  \qquad
  \text{\met{skill_dim_i}} = 1 - \frac{\text{\met{mse_dim_i}}}{\text{\met{persist_dim_i}}},
\end{equation}
%
together with \met{persist_loss}, the cap-scaled loss of the echo itself.
The \emph{skill score} is positive only when the model earns its parameters.
It is the stricter of the two ratios: a run can show a comfortable
\met{r2_dim_i} and a negative \met{skill_dim_i} at the same time, which means
it has learnt to reproduce the current velocity and little else.

\subsection{Amplitude and Restart Ratio}
\label{sec:metrics-amplitude}

A squared error cannot see one failure that matters a great deal for a
deployed controller. A model trained by least squares regresses towards the
conditional mean of the demonstration: where the pilot sometimes commanded a
strong manoeuvre and sometimes none, the model learns to command a weak one
every time. This is not a bug of the optimiser---it is what minimising
squared error does under uncertainty---but a robot flown by such a model is
sluggish, and the squared error \emph{rewards} the behaviour rather than
exposing it. Two metrics exist to expose it.

\paragraph{Amplitude ratio.}
The trainer logs the first and second moments of both prediction and target
per component (\met{sq_pred_dim_i} $= \mathbb{E}[\hat{y}_i^{2}]$,
\met{mean_pred_dim_i} $= \mathbb{E}[\hat{y}_i]$ and
\met{mean_true_dim_i} $= \mathbb{E}[y_i]$; the second moment of the target is
\met{baseline_dim_i}) and forms, once per epoch,
%
\begin{equation}
  \text{\met{amplitude_dim_i}}
  = \sqrt{\frac{\operatorname{Var}(\hat{y}_i)}{\operatorname{Var}(y_i)}}
  = \sqrt{\frac{\mathbb{E}[\hat{y}_i^{2}] - \mathbb{E}[\hat{y}_i]^{2}}
               {\mathbb{E}[y_i^{2}] - \mathbb{E}[y_i]^{2}}}.
\end{equation}
%
A value of one means the model reproduces the spread of the demonstrated
commands; a value well below one means it returns something closer to their
average than to any of them. The variance rather than the raw second moment
is used because on a component with a large mean the second moment is
dominated by that mean, and the spread---which is what collapses when a head
averages two modes---barely moves it.

\paragraph{Restart ratio.}
The frames that decide whether a deployed policy ever leaves standstill are
those in which the robot is at rest and the demonstration starts moving. They
are a fraction of the corpus small enough to be invisible in any aggregate, so
they are isolated explicitly. A frame is a \emph{restart frame} when the
measured velocity is below \SI{0.02}{} on every component and the demonstrated
command exceeds \SI{0.05}{} on at least one (in each component's own units).
Over those frames the trainer logs their count, \met{restart_frames}, and
%
\begin{gather}
  \text{\met{restart_pred_dim_i}} = \mathbb{E}\!\left[\hat{y}_i \cdot \operatorname{sign}(y_i)\right],
  \qquad
  \text{\met{restart_true_dim_i}} = \mathbb{E}\!\left[\,|y_i|\,\right],
  \\
  \text{\met{restart_ratio_dim_i}} = \frac{\text{\met{restart_pred_dim_i}}}{\text{\met{restart_true_dim_i}}}.
\end{gather}
%
The prediction is projected onto the demonstrated direction rather than taken
in absolute value: an absolute value would count noise around zero as if it
were a command and read high for a model that never moves. A restart ratio
near one means the model commands the departure the pilot commanded; near zero
it means the model would leave the robot sitting at its spawn. As with the
other ratios, the two halves are logged separately and divided once per epoch.

\subsection{Still and Moving Frames}
\label{sec:metrics-still}

Around \SI{35}{\percent} of the samples in the corpus are all-zero commands.
These are legitimate ``hold still'' targets---the pilot waiting, hovering, or
paused---and not missing labels, but they are also the easiest samples to
score well on. \met{mse_still} and \met{mse_moving} split the plain squared
error by whether the target is exactly zero on every component, so that it is
possible to tell how much of a score comes from recognising a stop and how
much from reproducing a manoeuvre. A model that scores well overall but poorly
on \met{mse_moving} is scoring on the stops.

Both columns are emitted for every batch, as \emph{not a number} when the
batch contains no frame of that kind; the epoch average is then taken only
over the batches in which the quantity was defined, so that the reported
value means what its name says.

\subsection{Per-Robot Breakdown}
\label{sec:metrics-per-robot}

The corpus is recorded on more than one platform, and a drone and a ground
robot do not use the same command components: vertical velocity is not merely
rare for a ground robot, it is unreachable. A fleet average can therefore
hide one robot being served well and another badly. Whenever the batch
carries robot identifiers, every per-dimension pair is repeated per model as
\met{mse_m<id>_dim_i} and \met{baseline_m<id>_dim_i}, and the logger derives
\met{r2_m<id>_dim_i} from them with no further bookkeeping. Because a batch
never straddles two recordings when the point cloud is in play, each batch
holds exactly one robot model, and the per-model averages are taken only over
the batches in which that model appears.

\subsection{Calibrated Metrics}
\label{sec:metrics-calibrated}

The regression-to-the-mean effect described in
Section~\ref{sec:metrics-amplitude} is corrected before deployment by an
\emph{output calibration}: a per-component affine map stored inside the
weights and applied by the inference node without knowing it exists,
%
\begin{equation}
  \tilde{y}_i = g_i\,\hat{y}_i + b_i,
  \qquad
  g_i = \frac{\operatorname{Cov}(\hat{y}_i, y_i)}{\operatorname{Var}(\hat{y}_i)},
  \qquad
  b_i = \mathbb{E}[y_i] - g_i\,\mathbb{E}[\hat{y}_i],
\end{equation}
%
which is the least-squares regression of the demonstrated command on the
model's prediction. The gain is fitted on the \emph{validation} split, not on
the training split: a model is already self-calibrated where it was
optimised, so a fit made there would read a gain near one and under-correct
everywhere else. It is refitted once per epoch so that every saved checkpoint
carries the calibration belonging to its own weights, and a component on
which either side carries no signal is left switched off ($g_i = b_i = 0$),
since a demonstration that never commands it has nothing to fit.

Because the calibrated output is what the robot is actually given, the whole
family of validation metrics is computed a second time through the fitted
calibration and logged with the prefix \met{val_cal_}: \met{val_cal_loss},
\met{val_cal_r2_dim_i}, \met{val_cal_amplitude_dim_i},
\met{val_cal_restart_ratio_dim_i}, and so on. The uncalibrated
\met{val_} family describes a model that is never deployed; the
\met{val_cal_} family describes the one that is.

The calibrated numbers are mildly optimistic, since the correction was fitted
on the same frames it is scored on. It is two parameters per component
against a whole split, so the optimism is small, but it is the reason
\met{val_loss}---fitted to nothing---remains the quantity that ranks
checkpoints.

\subsection{Model Health}
\label{sec:metrics-health}

The metrics above describe the output. A further group describes whether the
parts of the network that produce it are doing anything, which no loss curve
reveals.

\paragraph{Encoder spread.}
The architectures that consume a point cloud pass it through a sparse
convolutional encoder that reduces tens of thousands of voxels to a
200-dimensional descriptor (Section~\ref{sec:sparse-encoder}). \met{encoder/mean} and
\met{encoder/norm} are the mean value and the mean Euclidean norm of that
descriptor over the batch, and \met{encoder/std} is its standard deviation
\emph{across the samples of the batch}, per feature, averaged over features.
A value near zero means the encoder answers nearly the same vector whatever
cloud it is shown: the sparse branch is contributing nothing, the policy is
effectively flying on odometry alone, and the loss curve shows nothing of it.

The across-sample form is the only one that can detect this. The encoder ends
with a layer normalisation that pins every single descriptor to zero mean and
unit deviation across its channels, so a spread measured \emph{inside} one
descriptor reads $1.0$ whatever the input and the check could never fire.
Measured on this project, a trained encoder reports \met{encoder/std} around
$0.99$, while a randomly initialised one reports $0.08$. When a batch holds a
single sample the column is omitted, because $0.0$ would read as the failure
rather than as the absence of a measurement.

\paragraph{Hidden-state norms.}
For every recurrent cell $k$ in the network, \met{state/cell_<k>/mean},
\met{state/cell_<k>/std} and \met{state/cell_<k>/max} are the mean,
standard deviation and maximum over the batch of the Euclidean norm of the
cell's hidden state after the forward pass. They catch the two ways a
recurrent state goes wrong: growth without bound, and collapse to zero.

\paragraph{Gradient norms.}
\met{grad/total} is the global $L_2$ norm of the gradient before clipping,
the most immediate signal that training has stalled, and it is the one extra
number printed on the console at every epoch. \met{grad/<module>} is the same
norm restricted to each top-level sub-module of the network (the encoder, the
input layer, each cell, the output head), which is where a dead branch shows
up: a module absent from the list received no gradient at all in that batch,
which is not the same as receiving a small one, and the distinction separates
an unused branch from a merely quiet one. Gradient norms exist only for the
training split, since evaluation has no backward pass.

\paragraph{Cost.}
\met{epoch_time_s} is the wall-clock duration of the epoch, which together
with the number of parameters and the number of samples gives the throughput
of each architecture.

\subsection{Summary}

Table~\ref{tab:metrics-supervised} lists the supervised metrics by the
question they answer. Each is logged for the training and validation splits,
and the validation ones additionally through the calibration.

\begin{table}[htbp]
  \centering
  \small
  \caption{Supervised training metrics, grouped by the question they answer.
  The \texttt{train\_}, \texttt{val\_} and \texttt{val\_cal\_} prefixes are
  omitted. Values marked $\star$ are derived once per epoch from the logged
  pair rather than averaged over batches.}
  \label{tab:metrics-supervised}
  \begin{tabular}{@{}>{\raggedright\arraybackslash}p{0.34\textwidth}>{\raggedright\arraybackslash}p{0.60\textwidth}@{}}
    \toprule
    \textbf{Metric} & \textbf{What it measures} \\
    \midrule
    \multicolumn{2}{@{}l}{\emph{Headline}} \\
    \met{loss} & cap-scaled MSE, Eq.~\eqref{eq:capscaled}; \met{val_loss} ranks checkpoints \\
    \met{lr}, \met{epoch_time_s} & learning rate in effect; wall-clock cost of the epoch \\
    \addlinespace
    \multicolumn{2}{@{}l}{\emph{Is it learning anything?}} \\
    \met{baseline_zero_loss} & the loss of always predicting zero \\
    \met{mse_dim_i}, \met{baseline_dim_i} & per-component error and zero-predictor energy \\
    \met{r2_dim_i}\,$\star$ & $1 - \text{mse}/\text{baseline}$: 0 is no better than zero, 1 is perfect \\
    \met{persist_loss}, \met{persist_dim_i} & error of the velocity-echo predictor \\
    \met{skill_dim_i}\,$\star$ & $1 - \text{mse}/\text{persist}$: positive only when the model beats the echo \\
    \addlinespace
    \multicolumn{2}{@{}l}{\emph{What behaviour is it reproducing?}} \\
    \met{amplitude_dim_i}\,$\star$ & ratio of predicted to demonstrated spread; $<1$ is regression to the mean \\
    \met{restart_frames}, \met{restart_ratio_dim_i}\,$\star$ & how large the commanded departure is, on frames leaving standstill \\
    \met{mse_still}, \met{mse_moving} & error on hold-still targets against error on manoeuvres \\
    \met{r2_m<id>_dim_i}\,$\star$ & the per-component $R^2$ for each robot model in the fleet \\
    \addlinespace
    \multicolumn{2}{@{}l}{\emph{Is the machinery healthy?}} \\
    \met{encoder/std}, \met{encoder/mean}, \met{encoder/norm} & across-sample spread of the point-cloud descriptor; near zero means a dead encoder \\
    \met{state/cell_<k>/mean,std,max} & hidden-state norms per recurrent cell \\
    \met{grad/total}, \met{grad/<module>} & gradient norm, globally and per sub-module \\
    \bottomrule
  \end{tabular}
\end{table}

\section{Reinforcement Learning Metrics}
\label{sec:metrics-rl}

The reinforcement trainer takes the network the imitation stage produced and
refines it by flying: same architecture, same weights, now shaped by the
consequences of its own actions rather than by a pilot's example. The
algorithm is a recurrent variant of Proximal Policy Optimization
\cite{schulman2017ppo} with Generalised Advantage Estimation
\cite{schulman2016gae}, described in Section~\ref{sec:rppo-gae}. Its metrics
are logged once per iteration and fall into six groups: what the policy
achieved on the task; how the policy update behaved; how well the critic is
doing; what the policy's action distribution looks like; what the reward was
made of and what the vehicle flew through; and what the run cost. The first
group is the one a reader cares about; the others exist because, at one point
or another during this work, a run failed in a way the first group could not
explain.

\subsection{Task Performance}
\label{sec:metrics-rl-task}

\paragraph{Episodic return and length.}
For every episode that closed during the iteration, the trainer records the
undiscounted sum of the environment's rewards and the number of steps. They
are reported as \met{ep_return_mean}, \met{ep_return_std},
\met{ep_return_min}, \met{ep_return_max}, \met{ep_length_mean} and
\met{ep_length_std}, together with \met{n_episodes}, the number of episodes
those statistics average, and \met{episodes_total}, the running count since
the start of the run. The distinction between the last two matters:
\met{n_episodes} at zero says that nothing closed \emph{just now}, which is
normal on a fleet whose episodes last a thousand steps, whereas
\met{episodes_total} at zero says that nothing has \emph{ever} closed, which
is a fault. The return is computed from the environment's own reward, before
the truncation bootstrap described in Section~\ref{sec:rppo-gae} is added:
that bootstrap is right for the advantage trace and wrong for a reported
return, which would otherwise move with the critic on every truncated step.

\paragraph{Success rate and end reasons.}
The return alone cannot tell an episode that arrived from one that crashed.
Every episode ends for a named reason, and the trainer reports the fraction
of the iteration's closed episodes that ended for each as \met{term/<reason>}.
On the fleet the reasons are \texttt{goal} (arrived within tolerance),
\texttt{collision}, \texttt{flipped} (attitude beyond the terminal angle),
\met{out_of_band} (left the altitude envelope), \met{sensor_lost} (the
lidar returned nothing for two consecutive steps) and \met{max_steps}
(the time budget expired). \met{success_rate} is the share of the first.
The set of columns is counted rather than declared: whatever the environment
names becomes a column, so a fleet that adds a terminal condition gets a
column without the trainer changing. A blank cell therefore means ``did not
occur'' when episodes closed and ``not measured'' when none did, and the
plotting code reads the shares against \met{n_episodes} for that reason.

Why this is the primary metric, and not the return, is best shown by a
measurement made during this work: on one run \met{term/flipped} went from
\SI{73}{\percent} to \SI{97}{\percent} of episodes while the return improved
from $-19$ to $-4.6$. With a dense progress term and a fixed crash penalty,
``get closer, then crash'' scores better than ``crash early'' and keeps
improving for as long as the agent gets further, so the return rises the
whole way and the task is never solved.

\paragraph{Progress closed per episode.}
\met{ep_progress_mean} and \met{ep_progress_std} are the metres closed
between an episode's first observation and its last, measured as the
difference in distance to the goal. The column exists because the return can
disagree with it, and when they disagree it is the return that misleads. The
progress term of the reward is a potential-based shaping
\cite{ng1999policy} whose discounted sum over an episode telescopes to
$d(s_0) - \gamma^{T} d(s_T)$; the second half does not vanish at a terminal,
so an episode that ends far from the goal still collects
$d(s_0)\,(1 - \gamma^{T})$---measured on the fleet, $+5.6$ on a short slice
and $+24.5$ on the long one, against an arrival bonus of $20$. A run whose
episodes all truncate therefore shows a comfortably positive return while
nothing has been achieved. On one 650-iteration run this column fell from
\SI{1.26}{\metre} to \SI{0.04}{\metre} while the return improved throughout.
It is absent, rather than zero, on an environment that reports no distance.

\paragraph{Episodes in flight.}
Three columns describe the episodes that are \emph{still running} at the end
of the rollout: \met{inflight_steps}, the age of the oldest open episode;
\met{inflight_budget}, the mean step budget of the open episodes as the
environment sized them; and \met{inflight_d_goal}, their mean current distance
to the goal. Without them a run in which nothing has finished yet reports
\emph{not a number} for every episodic column and nothing else, and a long
episode cannot be told from one that can never end. \met{inflight_budget}
holding a single value for the whole run is also how a fleet made of identical
slices reveals itself.

\subsection{Policy Optimisation Diagnostics}
\label{sec:metrics-rl-policy}

PPO improves the policy by maximising a clipped surrogate objective over
several passes on the same rollout. The quantities below are computed on every
mini-batch of the update and averaged over the iteration, except where
stated; in the recurrent case only the real, non-padded positions of a
sequence contribute.

\paragraph{Losses.}
\met{pg_loss} is the clipped policy-gradient surrogate,
$-\mathbb{E}[\min(r_t A_t,\ \operatorname{clip}(r_t, 1-\epsilon, 1+\epsilon)\,A_t)]$
with $r_t = \pi_\theta(a_t \mid s_t)/\pi_{\theta_{\text{old}}}(a_t \mid s_t)$
the probability ratio, where the advantages are normalised to zero mean and
unit variance within each mini-batch so that the step size does not depend on
the reward's scale. \met{entropy} is the mean entropy of the action
distribution, and \met{actor_loss} is the total actor objective,
$\text{\met{pg_loss}} - c_{\text{ent}}\cdot\text{\met{entropy}}$. When the environment
declares which action components the robot actually commands, both the
log-probability and the entropy are summed over the commanded components
only; on this fleet two of the six Twist components (roll and pitch rate) are
never commanded, and without the mask a third of the trust region was spent
on noise in them.

\paragraph{Approximate KL divergence and the trust region.}
\met{approx_kl} is the estimated Kullback--Leibler divergence between the
policy that collected the rollout and the one being updated, using the
low-variance estimator $\mathbb{E}[(r_t - 1) - \log r_t]$. It is the
quantity that decides when the update stops: after every mini-batch the
running mean of the estimate over the current epoch is compared to
\met{target_kl}, and when it exceeds the threshold the actor stops for the
rest of the update, since the data no longer describes the current policy.
\met{kl_early_stop} records whether that happened (0 or 1) and
\met{epochs_run} how many of the \met{n_epochs} passes the actor completed.
The critic keeps fitting either way, as it has no ratio and nothing to be
wrong about.

The estimator is heavy-tailed on the small mini-batches a recurrent update
uses---sixteen samples on the fleet's preset---so a single sample above the
threshold is more often an outlier than a movement. The stop therefore acts
on the epoch's running mean, read only once it has at least eight informative
estimates, and only mini-batches scored \emph{after} the first optimiser step
count as informative (before it, the ratio is exactly one and the estimate
exactly zero). \met{approx_kl_max} is the largest per-mini-batch estimate of
the iteration, the one series a mean would destroy. It is read beside
\met{approx_kl}: close together, the policy genuinely steps that far and the
levers are \met{target_kl} or the learning rate; far apart, the estimator is
noisy and the mini-batch is too small.

\paragraph{Clip fraction.}
\met{clip_frac} is the fraction of samples whose ratio left the interval
$[1-\epsilon, 1+\epsilon]$, i.e.\ the share of the batch that hit the hinge of
the surrogate. Near zero the policy is barely moving---steps too small, or
advantages near zero; very high, it is being pushed hard against the hinge on
every batch. \met{clip_eps} is the value of $\epsilon$ in effect, logged
because it is annealed over the run.

\paragraph{Replay consistency.}
\met{replay_logratio_max} is a correctness check specific to the recurrent
implementation. On the first mini-batch of the first epoch, before anything
has been optimised, the update's forward pass should reproduce the rollout's
log-probabilities exactly, so the log-ratio must be zero to floating-point
error (about $10^{-6}$). A larger value means the update's forward pass is not
the rollout's---almost always a hidden state that was restored incorrectly---
and its only other symptom is a \met{clip_frac} that reads like an aggressive
learning rate. Cheap to compute and included for that reason
\cite{engstrom2020implementation}.

\paragraph{Optimisation state.}
\met{lr_actor} and \met{lr_critic} are the two learning rates in effect,
logged separately because the actor and the critic have separate optimisers
and separate annealing. \met{actor_grad_norm} is the actor's gradient norm
before clipping, averaged over the optimiser steps actually taken (it is
undefined on mini-batches the early stop skipped, and those are excluded from
the mean rather than dragging it to zero).

\subsection{Value Function}
\label{sec:metrics-rl-value}

The critic estimates $V(s)$, and everything in the advantage rests on it. Its
health is described by the following columns.

\paragraph{Losses.}
\met{vf_loss} is the mean squared regression error of $V(s)$ onto the GAE
returns, optionally clipped so that the value cannot move more than a fixed
range from what it predicted during collection; \met{critic_loss} is the same
quantity times its coefficient. \met{critic_grad_norm} is the critic's
gradient norm before clipping.

\paragraph{Explained variance.}
The single most informative number about the critic is
%
\begin{equation}
  \text{\met{explained_var}}
  = 1 - \frac{\operatorname{Var}(R_t - V(s_t))}{\operatorname{Var}(R_t)},
\end{equation}
%
computed over the \emph{whole rollout}, with the values the rollout was
collected with. A value of $1$ is a perfect critic, $0$ one no better than
predicting the mean return, negative one worse than that.
\met{explained_var_batch} asks the same question of each mini-batch and
averages the answers, and it is logged beside the first only as a warning: a
mini-batch is a handful of sequences of \emph{consecutive} steps whose returns
barely vary, so its denominator is tiny and individual batches reach large
negative values that the mean inherits. Measured on a fleet run,
\met{explained_var_batch} sat near $-1.4$ for 390 iterations with a minimum of
$-58$, while the value estimate tracked the returns in mean and spread
throughout. The per-batch figure had already been acted on---the discount had
been cut on its evidence---before the discrepancy was understood, which is why
the rollout-level one is the number the diagnostic tables in
Chapter~\ref{ch:training} are calibrated against.

\paragraph{Value and return statistics.}
\met{value_mean}, \met{value_std}, \met{value_min} and \met{value_max}
describe the distribution of $V(s)$ over the update, and \met{return_mean}
and \met{return_std} that of the GAE returns it is regressed onto;
\met{value_mse} is their mean squared difference. \met{value_std} deserves
its own reading: exactly $0.0$ is not a small number, it is a constant
predictor. During this work a whole fleet run reported \met{value_std}
exactly $0.0$ and \met{explained_var} of $10^{-6}$ because the critic
received the observation on a different scale from the actor; every unit of
its first activation layer saturated, and the only thing the value head could
still learn was its output bias. The advantages were then simply the raw
returns.

\paragraph{Advantages.}
\met{adv_mean} and \met{adv_std} are the mean and spread of the advantages as
GAE produced them, \emph{before} the per-mini-batch normalisation. A mean far
from zero says that the critic and the returns have drifted apart, which the
losses alone do not distinguish from a hard task; a collapsing spread says
that there is nothing left to learn from, or that the reward has stopped
varying.

\subsection{Exploration and Action Statistics}
\label{sec:metrics-rl-actions}

For a continuous action space the policy is a Gaussian whose mean the network
produces and whose standard deviation is a trained parameter of the agent.
Five columns describe what that distribution is doing. With the action mask
on, all of them cover the commanded components only, which is a change of
\emph{units} rather than of behaviour when comparing against a run made
without it.

\begin{itemize}
  \item \met{action_std_mean} --- the mean standard deviation of the Gaussian.
    It is trained, so watching it is how one sees exploration collapsing
    (approaching zero: the entropy bonus is too low and the policy has stopped
    exploring) or running away.
  \item \met{action_abs_mean} --- the mean absolute value of the \emph{sampled}
    actions, i.e.\ how large the commands sent to the environment are.
  \item \met{action_sat_frac} --- the fraction of sampled action components that
    the environment's clamp altered. The buffer stores the unclamped sample
    and its log-probability, so a rising fraction means the update is learning
    about commands the robot never applied.
  \item \met{action_mu_abs_mean} --- the mean absolute value of the Gaussian's
    \emph{mean}, as opposed to its samples.
  \item \met{action_mu_outside_frac} --- the fraction of components on which the
    Gaussian's mean itself lies outside the action envelope. This is a
    different fact from saturation: a mean well outside the limits is a policy
    the clamp has made deterministic in those components---the sampling noise
    no longer changes what the robot is sent---while the entropy bonus keeps
    paying for exploration that no longer exists. Nothing in the standard
    objective charges for the excess, which is why an explicit bound penalty
    was added to the actor loss; this column is what says whether it is
    holding.
\end{itemize}

\subsection{Reward Decomposition and Flight Signals}
\label{sec:metrics-rl-reward}

\paragraph{Per-step reward.}
\met{reward_step_mean} and \met{reward_step_std} are the mean and standard
deviation of the environment's reward over every step of the rollout, with
no assumptions about how episodes end.

\paragraph{Contribution of each term.}
The rewards used on the fleet are composed of named terms---\texttt{progress},
\texttt{time}, \texttt{safety}, \texttt{tilt}, \texttt{altitude},
\texttt{smoothness} and \texttt{terminal}---each of which knows its own
arithmetic and its share of the per-step budget (Section~\ref{sec:rl-reward}).
The environment accumulates what each term contributed, and the trainer writes
the mean per-step contribution of each as \met{reward/<term>}. The values sum
to \met{reward_step_mean} exactly, so a gap between the two means something is
contributing without being logged. This decomposition is what makes the
weights tunable against a measurement rather than against behaviour: a term
that has taken the reward over---a time cost grown into the whole signal, a
terminal penalty carrying \SI{90}{\percent} of the return, which is what one
configuration of the safe-progress reward was found to do---is a number
rather than an inference. From the shares of the end reasons and the mean
episode length, the dense part of the return can also be reconstructed as
$\text{dense} = \text{\met{ep_return_mean}} + R_{\text{crash}}\,P(\text{crash}) -
R_{\text{goal}}\,P(\text{goal})$; a return that improves while the dense part
worsens is termination bias, not learning.

\paragraph{What the fleet flew through.}
A cost of zero has two causes: the robot stayed clear, or the term never
fired. The reward columns answer the second; a set of tracked signals answers
the first. For every signal the fleet reports, the trainer logs
\met{signal/<name>_mean}, its mean over the rollout;
\met{signal/<name>_worst}, the end of its range that means trouble; and
\met{signal/<name>_coverage}, the fraction of steps on which it could be
measured at all. The direction of \emph{worst} is declared per signal, because
one aggregate cannot serve both: the minimum of a clearance, the maximum of an
attitude---the minimum tilt over a rollout says nothing. Coverage qualifies
the other two: a term whose signal was absent charged nothing, which on its
own reads as a term with nothing to complain about. The signals are chosen so
that every terminal condition the fleet can hit has one, so that a
termination that dominates a run can be traced to how close the fleet was
flying to the threshold behind it; Table~\ref{tab:metrics-signals} lists them.

\begin{table}[htbp]
  \centering
  \small
  \caption{Signals tracked over each rollout, with the direction reported as
  \texttt{\_worst}. Counters are reported by how much they grew, fleet-wide
  (\texttt{\_total}) and on the worst single slice (\texttt{\_worst\_slice}).}
  \label{tab:metrics-signals}
  \begin{tabular}{@{}ll>{\raggedright\arraybackslash}p{0.52\textwidth}@{}}
    \toprule
    \textbf{Signal} & \textbf{Worst} & \textbf{What it says} \\
    \midrule
    \met{min_obstacle_dist_nofloor} & min & how close to anything the fleet actually flew, ground excluded \\
    \met{height_agl} & min & whether the altitude band is being held \\
    \met{tilt_gt} & max & how close to the attitude limit, which ends most episodes \\
    \met{d_goal} & min & the closest approach, when nothing arrives \\
    \met{yaw_error} & max & heading error against the one the mission asked for \\
    \met{speed_gt} & max & whether the vehicle is moving at all \\
    \met{pipeline_latency_ms} & max & the age of the observation \\
    \met{index_drift} & max & drift of the ground-truth frame the reward reads \\
    \addlinespace
    \textbf{Counter} & & \\
    \met{dropped_clouds} & -- & clouds the perception chain produced and the bridge could not use \\
    \met{stale_clouds} & -- & clouds older than the step that asked for them \\
    \met{advance_steps} & -- & physics steps spent: the simulation work the rollout cost \\
    \bottomrule
  \end{tabular}
\end{table}

Beside the distributions, the slices' own counters are reported as
\met{signal/<name>_total} and \met{signal/<name>_worst_slice}, by how much
they grew over the rollout. The slices zero them on every episode reset, so a
rollout sees them saw-tooth and the growth is accumulated across the resets.
Two numbers, because nothing in the reward reads them and only the pair can
tell two situations apart: a third of the clouds dropped on one slice is a
broken slice training on a world it is barely seeing, while the same total
spread evenly is a fleet running too hot.

\subsection{Normalisers and Network Health}
\label{sec:metrics-rl-health}

\met{obs_rms_var_mean} is the mean variance the external observation
normaliser is dividing by, and \met{ret_rms_var} the variance of the discounted
return the reward normaliser is dividing by, when either is enabled. Both are
lifetime estimates with no discount, so what remains to read in them is a
\emph{drift}, which is the distribution genuinely moving: a jump that then
decays slowly used to mean a single outlier had permanently inflated the
scale, which is why each sample is now bounded before being folded in.

\met{encoder/mean}, \met{encoder/std} and \met{encoder/norm} are the same
statistics as in Section~\ref{sec:metrics-health}, but measured over every
time step of the rollout rather than over one batch, since a rollout collects
one time step at a time. \met{state/cell_<k>/mean,std,max} are the hidden-state
norms, as before.

\subsection{Efficiency and Thermal Guard}
\label{sec:metrics-rl-cost}

\met{total_steps} and \met{time_elapsed_s} place the iteration in the run;
\met{sps} is the throughput in environment steps per second. The wall-clock
time of an iteration is split into \met{collect_time_s} (the rollout as a
whole), \met{env_time_s} (the part spent inside the environment's
\texttt{step}, i.e.\ the simulator), \met{policy_time_s} (the actor's forward
passes during collection) and \met{update_time_s} (the optimisation). One
number for all of them sends the optimisation effort to the wrong place, and
on the fleet it did: the simulator, not the trainer, is the bottleneck.

Long runs on a workstation are paused when the hardware overheats.
\met{gpu_temp_c} and \met{cpu_temp_c} are the temperatures read at the end of
the iteration and \met{thermal_pause_s} the wall-clock time the guard has
taken so far, so that a slow iteration can be told from a paused one.

\subsection{Summary}

Table~\ref{tab:metrics-rl} lists the reinforcement learning metrics by group.

\begin{longtable}{@{}>{\raggedright\arraybackslash}p{0.40\textwidth}>{\raggedright\arraybackslash}p{0.54\textwidth}@{}}
  \caption{Reinforcement learning metrics, one row per iteration, grouped by
  the question they answer.}
  \label{tab:metrics-rl} \\
  \toprule
  \textbf{Metric} & \textbf{What it measures} \\
  \midrule
  \endfirsthead
  \multicolumn{2}{@{}l}{\emph{Table~\thetable\ (continued)}} \\
  \toprule
  \textbf{Metric} & \textbf{What it measures} \\
  \midrule
  \endhead
  \bottomrule
  \endlastfoot
    \multicolumn{2}{@{}l}{\emph{Task performance}} \\
    \met{ep_return_{mean,std,min,max}} & episodic return of the episodes closed this iteration \\
    \met{ep_length_{mean,std}} & their length in steps \\
    \met{n_episodes}, \met{episodes_total} & episodes closed now, and since the start \\
    \met{success_rate}, \met{term/<reason>} & share of episodes that arrived, and share per end reason \\
    \met{ep_progress_{mean,std}} & metres closed per episode; the return's honest counterpart \\
    \met{inflight_{steps,budget,d_goal}} & state of the episodes still open \\
    \addlinespace
    \multicolumn{2}{@{}l}{\emph{Policy update}} \\
    \met{pg_loss}, \met{entropy}, \met{actor_loss} & clipped surrogate, policy entropy, total actor objective \\
    \met{approx_kl}, \met{approx_kl_max} & policy movement: epoch mean, and the largest mini-batch \\
    \met{target_kl}, \met{kl_early_stop}, \met{epochs_run}, \met{n_epochs} & the trust region and whether it cut the update short \\
    \met{clip_frac}, \met{clip_eps} & share of samples on the hinge; the $\epsilon$ in effect \\
    \met{replay_logratio_max} & recurrent-state correctness check; must be $\approx 0$ \\
    \met{lr_actor}, \met{lr_critic}, \met{actor_grad_norm} & optimiser state \\
    \addlinespace
    \multicolumn{2}{@{}l}{\emph{Value function}} \\
    \met{vf_loss}, \met{critic_loss}, \met{critic_grad_norm} & critic regression loss and gradient \\
    \met{explained_var}, \met{explained_var_batch} & variance of the returns the critic accounts for; rollout-level and per-batch \\
    \met{value_{mean,std,min,max}}, \met{value_mse} & distribution of $V(s)$ and its error; \met{value_std} $= 0$ is a constant critic \\
    \met{return_{mean,std}}, \met{adv_{mean,std}} & GAE returns and raw advantages \\
    \addlinespace
    \multicolumn{2}{@{}l}{\emph{Action distribution}} \\
    \met{action_std_mean} & width of the Gaussian: exploration \\
    \met{action_abs_mean}, \met{action_sat_frac} & size of the samples; how often the clamp fired \\
    \met{action_mu_abs_mean}, \met{action_mu_outside_frac} & size of the mean; how often it lies outside the envelope \\
    \addlinespace
    \multicolumn{2}{@{}l}{\emph{Reward and environment}} \\
    \met{reward_step_{mean,std}} & the reward per step \\
    \met{reward/<term>} & mean per-step contribution of each term; sums to the above \\
    \met{signal/<name>_{worst,mean,coverage}} & what the fleet flew through \\
    \met{signal/<name>_{total,worst_slice}} & growth of the slices' own counters \\
    \addlinespace
    \multicolumn{2}{@{}l}{\emph{Health and cost}} \\
    \met{obs_rms_var_mean}, \met{ret_rms_var} & scales of the observation and return normalisers \\
    \met{encoder/*}, \met{state/cell_*} & encoder spread and hidden-state norms, over the rollout \\
    \met{sps}, \met{*_time_s} & throughput and where an iteration's time goes \\
    \met{gpu_temp_c}, \met{cpu_temp_c}, \met{thermal_pause_s} & temperatures and time lost to the thermal guard \\
\end{longtable}
